\chapter*{Conventions used in code and examples}
\addcontentsline{toc}{chapter}{Conventions used in code and examples}

Code is printed in a monospaced typeface. A command appears as
\cmd{section}; a package as \pkg{graphicx}; and a file as \file{main.tex}.
Required material is written in braces, as in \required{article}. Optional
material is written in brackets, as in \optional{11pt}. A descriptive item
that you must replace appears as \placeholder{your title}.

Complete code panels can be copied into a plain-text source file. Their
quotation marks and commands use straight ASCII characters. A code panel that
is intentionally invalid is labelled \emph{Broken example: do not copy as a
working document}. The corresponding repair follows nearby.

The recurring boxes have consistent jobs:

\begin{description}[style=nextline,leftmargin=1.6em]
  \item[Classroom Note] clarifies an idea without interrupting the main path.
  \item[Try It Now] asks for an immediate, small action.
  \item[Pause and Predict] asks you to decide what will happen before compiling.
  \item[Common Error] shows a realistic failure, cause, and repair.
  \item[Good Practice] distinguishes maintainable professional habits.
  \item[Science in Practice] connects a command to scholarly work.
  \item[Advanced Topic] marks material that a beginner may safely postpone.
  \item[Submission Check] turns a requirement into an auditable check.
  \item[Mini Project] develops a cumulative document.
\end{description}

Values in the continuing article are labelled as simulated demonstration
data. The examples do not claim genuine scientific results. Web interfaces
may change; descriptive instructions identify the task rather than depending
on an exact button position.
